\documentclass[conference]{IEEEtran}

\usepackage{amsmath,amssymb}
\usepackage{textcomp}

\begin{document}

\title{SGN Loss Function Calculation}

\author{\IEEEauthorblockN{Raymond Tie Siou Hin}
\IEEEauthorblockA{\textit{Monash University Malaysia}\\
Bandar Sunway, Malaysia}
}

\maketitle

\begin{abstract}
This note summarizes how the loss function is calculated in the original
Subtask Gated Network (SGN) formulation and how the current implementation
extends it to the multi-appliance case.
\end{abstract}

\begin{IEEEkeywords}
non-intrusive load monitoring, SGN, loss function, appliance classification,
power regression
\end{IEEEkeywords}

\section{Original SGN Paper Loss}

In the original SGN paper, the model is formulated for one target appliance
$i$ at a time. For an output sequence of length $T$, the ground-truth power
sequence is
\begin{equation}
\mathbf{y}^{i} = \{y^{i}_{1}, y^{i}_{2}, \ldots, y^{i}_{T}\}.
\end{equation}

The on/off label is generated using a 15 W threshold:
\begin{equation}
o^{i}_{t} =
\begin{cases}
1, & y^{i}_{t} > 15\text{ W},\\
0, & y^{i}_{t} \leq 15\text{ W}.
\end{cases}
\end{equation}

SGN has two subnetworks. The regression subnetwork predicts appliance power:
\begin{equation}
\hat{\mathbf{p}}^{i}
=
f^{i}_{power}(\tilde{\mathbf{x}}),
\end{equation}
and the classification subnetwork predicts the appliance on probability:
\begin{equation}
\hat{\mathbf{o}}^{i}
=
f^{i}_{on}(\tilde{\mathbf{x}}).
\end{equation}

The final SGN output is the element-wise product of the two subnetworks:
\begin{equation}
\hat{\mathbf{y}}^{i}
=
\hat{\mathbf{p}}^{i} \odot \hat{\mathbf{o}}^{i},
\end{equation}
where $\odot$ denotes element-wise multiplication.

The gated output regression loss is
\begin{equation}
\mathcal{L}^{i}_{output}
=
\frac{1}{T}
\sum_{t=1}^{T}
\left(
y^{i}_{t}
-
\hat{p}^{i}_{t}\hat{o}^{i}_{t}
\right)^2.
\end{equation}

The on/off classification loss is binary cross entropy:
\begin{equation}
\mathcal{L}^{i}_{on}
=
-
\frac{1}{T}
\sum_{t=1}^{T}
\left[
o^{i}_{t}\log(\hat{o}^{i}_{t})
+
(1-o^{i}_{t})\log(1-\hat{o}^{i}_{t})
\right].
\end{equation}

The total SGN loss for appliance $i$ is
\begin{equation}
\mathcal{L}^{i}
=
\mathcal{L}^{i}_{output}
+
\mathcal{L}^{i}_{on}.
\end{equation}

Therefore, the original SGN loss is appliance-indexed. The notation
$\mathcal{L}^{i}$ means the loss is calculated for one appliance, such as
fridge or microwave. The original formulation does not describe one shared
multi-output loss averaged across all appliances.

\section{Current Multi-Appliance Implementation Loss}

The current implementation can train one model with multiple appliance outputs.
The model tensors have shape
\begin{equation}
[B, A, T],
\end{equation}
where $B$ is batch size, $A$ is the number of target appliances, and $T$ is
the output sequence length.

For batch item $b$, appliance $a$, and timestep $t$, the model predicts
\begin{equation}
\hat{P}_{b,a,t}
\quad\text{and}\quad
\hat{O}_{b,a,t},
\end{equation}
where $\hat{P}_{b,a,t}$ is the regression power prediction and
$\hat{O}_{b,a,t}$ is the on/off probability.

The gated prediction is
\begin{equation}
\hat{Y}_{b,a,t}
=
\hat{P}_{b,a,t}\hat{O}_{b,a,t}.
\end{equation}

The implementation computes the output error as
\begin{equation}
E^{output}_{b,a,t}
=
\left(
Y_{b,a,t}
-
\hat{Y}_{b,a,t}
\right)^2,
\end{equation}
and the on/off classification error as
\begin{equation}
E^{on}_{b,a,t}
=
-
\left[
O_{b,a,t}\log(\hat{O}_{b,a,t})
+
(1-O_{b,a,t})\log(1-\hat{O}_{b,a,t})
\right].
\end{equation}

The overall output loss is averaged across batch, appliances, and time:
\begin{equation}
\mathcal{L}_{output}
=
\frac{1}{BAT}
\sum_{b=1}^{B}
\sum_{a=1}^{A}
\sum_{t=1}^{T}
E^{output}_{b,a,t}.
\end{equation}

The overall classification loss is also averaged across batch, appliances,
and time:
\begin{equation}
\mathcal{L}_{on}
=
\frac{1}{BAT}
\sum_{b=1}^{B}
\sum_{a=1}^{A}
\sum_{t=1}^{T}
E^{on}_{b,a,t}.
\end{equation}

The final implementation loss is
\begin{equation}
\mathcal{L}
=
\mathcal{L}_{output}
+
\mathcal{L}_{on}.
\end{equation}

This keeps the SGN idea of combining gated power regression with on/off
classification. However, when $A>1$, the implementation averages the loss
over all selected appliances. This differs from the original paper's
appliance-indexed formulation.

\section{Code Mapping}

In the implementation, the loss is calculated in
\texttt{baseline/SGN/sgn/losses.py}. The main operations are:
\begin{equation}
\texttt{output\_error}
=
\left(
\texttt{gated\_power}
-
\texttt{target\_power}
\right)^2,
\end{equation}
\begin{equation}
\texttt{on\_error}
=
\texttt{BCE}(
\texttt{on\_prob},
\texttt{target\_on}
),
\end{equation}
\begin{equation}
\texttt{loss}
=
\texttt{mean(output\_error)}
+
\texttt{mean(on\_error)}.
\end{equation}

For one appliance, this is equivalent to the original SGN loss. For multiple
appliances, this becomes a shared average loss across all selected appliances.

\section{Numerical Example}

Assume one appliance and an output length of three timesteps:
\begin{equation}
\mathbf{y} = [0, 100, 120],
\qquad
\mathbf{o} = [0, 1, 1].
\end{equation}
Suppose the regression branch predicts
\begin{equation}
\hat{\mathbf{p}} = [20, 90, 130],
\end{equation}
and the classification branch predicts
\begin{equation}
\hat{\mathbf{o}} = [0.1, 0.8, 0.6].
\end{equation}

The gated output is
\begin{equation}
\hat{\mathbf{y}}
=
\hat{\mathbf{p}}\odot\hat{\mathbf{o}}
=
[2, 72, 78].
\end{equation}

The output loss is
\begin{equation}
\mathcal{L}_{output}
=
\frac{(0-2)^2 + (100-72)^2 + (120-78)^2}{3}
=
850.67.
\end{equation}

The on/off loss is
\begin{equation}
\mathcal{L}_{on}
=
-
\frac{
\log(0.9)
+
\log(0.8)
+
\log(0.6)
}{3}
=
0.280.
\end{equation}

The total loss is
\begin{equation}
\mathcal{L}
=
850.67 + 0.280.
\end{equation}

In the actual implementation, appliance power is normalized before training,
so the numerical regression loss is much smaller than this raw-watt example.

\section{Practical Reproduction Note}

To reproduce the original SGN formulation more closely, train one appliance at
a time. The multi-appliance version is useful for a shared NILM pipeline, but
it is not the exact appliance-indexed setup described in the original SGN
paper.

\end{document}
